\documentclass[11pt,a4paper]{article}

% Пакеты
\usepackage[utf8]{inputenc}
\usepackage[T2A]{fontenc}
\usepackage[russian]{babel}
\usepackage[left=1.5cm,right=1.5cm,top=1.5cm,bottom=1.5cm]{geometry}
\usepackage{graphicx}
\usepackage{xcolor}
\usepackage{titlesec}
\usepackage{enumitem}
\usepackage{hyperref}
\usepackage{multicol}
\usepackage{fontawesome}

% Цветовая схема
\definecolor{primarycolor}{RGB}{14, 94, 157}
\definecolor{secondarycolor}{RGB}{102, 102, 102}
\definecolor{accentcolor}{RGB}{231, 76, 60}
\definecolor{pastelblue}{RGB}{80, 140, 200}
\definecolor{pastelgreen}{RGB}{70, 150, 120}
\definecolor{pastelpurple}{RGB}{120, 100, 160}
\definecolor{pastelteal}{RGB}{60, 140, 150}
\definecolor{pastelorange}{RGB}{180, 130, 90}
\definecolor{pastelrose}{RGB}{160, 100, 120}
\definecolor{pasteldot}{RGB}{160, 190, 220}

\hypersetup{
    colorlinks=true,
    linkcolor=primarycolor,
    urlcolor=primarycolor,
    pdfborder={0 0 0}
}

\setlength{\parindent}{0pt}
\setlength{\parskip}{0.5em}

\titleformat{\section}
  {\Large\bfseries\color{primarycolor}}
  {}{0em}{}[\titlerule]
\titlespacing*{\section}{0pt}{12pt}{6pt}

\titleformat{\subsection}
  {\large\bfseries\color{secondarycolor}}
  {}{0em}{}
\titlespacing*{\subsection}{0pt}{8pt}{4pt}

\setlist[itemize]{leftmargin=*, itemsep=2pt, parsep=0pt, topsep=4pt, label={\color{pasteldot}\textbullet}}

\newcommand{\hlblue}[1]{\textcolor{pastelblue}{#1}}
\newcommand{\hlgreen}[1]{\textcolor{pastelgreen}{#1}}
\newcommand{\hlpurple}[1]{\textcolor{pastelpurple}{#1}}
\newcommand{\hlteal}[1]{\textcolor{pastelteal}{#1}}
\newcommand{\hlorange}[1]{\textcolor{pastelorange}{#1}}
\newcommand{\hlrose}[1]{\textcolor{pastelrose}{#1}}

\newcommand{\name}[1]{
  \begin{center}
    {\Huge\bfseries\color{primarycolor} #1}
  \end{center}
  \vspace{-8pt}
}

\newcommand{\tagline}[1]{
  \begin{center}
    {\large\color{secondarycolor} #1}
  \end{center}
  \vspace{4pt}
}

\newcommand{\contactinfo}[4]{
  \begin{center}
    \color{secondarycolor}
    #1 \quad | \quad #2 \quad | \quad #3 \quad | \quad #4
  \end{center}
  \vspace{8pt}
}

\newcommand{\cventry}[5]{
  \noindent
  \textbf{\color{primarycolor}#1} \hfill {\color{secondarycolor}\textit{#2}}\\
  {\color{secondarycolor}#3} \hfill {\color{secondarycolor}#4}
  \vspace{2pt}

  #5
  \vspace{6pt}
}

\newcommand{\achievement}[4]{
  \noindent
  \textbf{\color{primarycolor}#1} \hfill {\color{secondarycolor}\textit{#2}}\\
  {\color{secondarycolor}#3}
  \vspace{2pt}

  #4
  \vspace{6pt}
}

\begin{document}

% Заголовок с фото
\begin{minipage}[t]{0.7\textwidth}
  \name{Никита Новицкий}
  \tagline{Студент направления «Математика и компьютерные науки»}
  \contactinfo
    {\small +7 (977) 414-09-73}
    {\tiny n.novitskiy@edu.centraluniversity.ru \normalsize}
    {Москва, Россия}
    {\small \href{https://github.com/xWooshieL}{github.com/xWooshieL}}
\end{minipage}
\hfill
\begin{minipage}[t]{0.25\textwidth}
  \raggedleft
  \IfFileExists{photo_3.5x4.5cm_300dpi.jpg}{%
    \includegraphics[width=3.5cm,height=3.5cm,keepaspectratio]{photo_3.5x4.5cm_300dpi.jpg}%
  }{}
\end{minipage}

\vspace{8pt}

\section{О себе}

Студент \hlblue{Центрального Университета}  на бюджетной основе по направлению «Математика и компьютерные науки». Глубоко интересуюсь \hlgreen{машинным обучением}, \hlpurple{искусственным интеллектом} и их применением в научных исследованиях. Обладаю сильными навыками в программировании на Python, анализе данных и статистике.

Целеустремлённый и усидчивый. При работе над проектами могу посвящать им до 16 часов в сутки, достигая глубокого понимания предмета. Стремлюсь развиваться в области AI/ML и внести вклад в инновационные проекты ведущих технологических компаний России.

\section{Образование}

\cventry
  {Центральный Университет}
  {2025-- настоящее время}
  {Бакалавриат: Математика и компьютерные науки}
  {Москва}
  {
    \begin{itemize}
      \item Обучение на \hlteal{бюджетной} основе
      \item Специализация: Machine Learning и Artificial Intelligence
      \item Математические курсы: основы линейной алгебры и математического анализа, дискретная математика
      \item Продвинутое программирование: Python
      \item Статистика и теория вероятностей, введение в искусственный интеллект, экономика
      \item Участие в \hlorange{Научной студии} по переменным звёздам (применение ML в астрономии)
      \item Участие в проекте \hlblue{MusicTech} (ИИ для музыкального образования, при поддержке Т-Банка)
    \end{itemize}
  }

\cventry
  {НИУ ВШЭ (Высшая Школа Экономики)}
  {2024 -- 2025}
  {Бакалавриат: ФКН ПМИ}
  {Москва}
  {
    \begin{itemize}
      \item Обучение на бюджетной основе
      \item Математические курсы: линейная алгебра и геометрия, математический анализ, дискретная математика, теория чисел, алгебра
      \item Программирование: введение в Python, Алгоритмы и структуры данных, введение с C++
    \end{itemize}
  }

\cventry
  {МАОУ СОШ №3}
  {2024}
  {Среднее общее образование}
  {г. Наро-Фоминск}
  {
    \begin{itemize}
      \item 10-11 классы с углублённым изучением английского и математики
    \end{itemize}
  }

\section{Достижения в период студенчества (2024--2026)}

\achievement
  {2 место · буткемп М.Видео (Query Understanding / NER)}
  {2026}
  {Промышленный кейс М.Видео по умному поиску и извлечению фактов из запросов}
  {
    \begin{itemize}
      \item Команда (Никита · Некит · Лиза) заняла \hlblue{2 место} по итогам буткемпа
      \item Гибридный каскад: SpellFix $\to$ словари/regex $\to$ CRF NER (BIO) $\to$ классификаторы бренда, категории и типов атрибутов
      \item Brand clf: 67 классов + NO\_BRAND/UNKNOWN; честный gold-eval на 823 запросах
      \item Эксперименты с GLiNER; тот же каскад в FastAPI/CLI и нативном Qt SmartSearch
      \item GitHub: \href{https://github.com/xWooshieL/mvideo-ner-search}{github.com/xWooshieL/mvideo-ner-search}
    \end{itemize}
  }

\achievement
  {Доклад на конференции «Научный Телеграф» · MusicTech}
  {2026}
  {Мультидисциплинарная молодёжная конференция Центрального университета}
  {
    \begin{itemize}
      \item Проект \hlgreen{MusicTech} (ЦУ $\times$ Т-Банк): виртуальный оркестр / real-time score following на ИИ
      \item Вместе с командой собрали MVP (HMM/HSMM/OLTW + RL для темпа), написали статью с математическими постановками
      \item Выступление 17 мая 2026: этапы разработки и гипотезы (конференция: 500+ участников, 200+ докладов)
      \item Статья: \href{https://github.com/xWooshieL/music-tech/blob/master/article/main.pdf}{main.pdf} ·
        слайды: \href{https://github.com/xWooshieL/music-tech/blob/master/article/presentation/presentation.pdf}{presentation.pdf} ·
        \href{https://musictech.tech}{musictech.tech}
    \end{itemize}
  }

\achievement
  {Разработка системы автоматического ценообразования ГК МОС}
  {2026}
  {Система поиска и верификации цен на строительные материалы для компании ГК МОС}
  {
    \begin{itemize}
      \item Модуль извлечения данных: SpaCy NER (дообучен), regex, OpenRouter LLM
      \item Модуль поиска цен: 5-уровневый каскад, Yandex API, Crawl4AI, BS4
      \item Модуль верификации TPA (модель верификации, которая проверяет, что найденный на сайте товар совпадает с позицией из спецификации): sentence-transformers, Soundex, Levenshtein, CatBoost, Spec Match по техническим характеристикам
      \item Модуль выбора цены: ранжирование, учёт НДС, отсечение выбросов
      \item База знаний из цен, быстрый поиск по TPA для повторяющихся позиций
      \item Python, 350+ unit-тестов
    \end{itemize}
  }

\achievement
  {Абсолютный победитель хакатона MGIMO GPT CHALLENGE}
  {2025}
  {Престижный хакатон по промпт инжиниирингу и искусственному интеллекту}
  {
    \begin{itemize}
      \item Прохождение в \hlblue{финальный этап} из сотен участников
      \item Первое место в финальном этапе и приз в размере 200.000 рублей
      \item Разработка инновационных решений в области NLP и Machine Learning
      \item Работа над сложными задачами анализа текстов и генерации
      \item Применение передовых методов prompt engineering и feature engineering
    \end{itemize}
  }

\achievement
  {Автор научной публикации: метод MLFE}
  {2025}
  {Multi-Level Feature Engineering в Prompt Engineering}
  {
    \begin{itemize}
      \item \hlgreen{Название:} «MULTI-LEVEL FEATURE ENGINEERING В PROMPT ENGINEERING: ПОВЫШЕНИЕ ЭФФЕКТИВНОСТИ ГЕНЕРАЦИИ ЧЕРЕЗ ИЕРАРХИЧЕСКИЙ СИНТЕЗ ПРИЗНАКОВ»
      \item Разработка нового метода повышения качества работы больших языковых моделей
      \item Научная статья опубликована в рамках хакатона MGIMO GPT CHALLENGE
      \item GitHub: \href{https://github.com/xWooshieL/interview_analysis_pipeline/blob/main/MLFE\%20Novitskiy.pdf}{github.com/xWooshieL/interview\_analysis\_pipeline}
    \end{itemize}
  }

\achievement
  {Участник научной студии по переменным звёздам}
  {2025}
  {Центральный Университет}
  {
    \begin{itemize}
      \item Применение методов машинного обучения для \hlpurple{предсказания переменности звёзд}
      \item Работа с реальными астрономическими данными и временными рядами
      \item Участие в Kaggle соревновании в рамках исследования научной студии университета
      \item Разработка, обучение и оптимизация собственных моделей классификации
      \item Использование математики и ML для решения научных задач
      \item Собственная статья с решением проблемы
      \item Обучали модель с использованием sklearn, CatBoost, xgboost и стандартных ML-алгоритмов
      \item Kaggle: \href{https://www.kaggle.com/competitions/cu-science-studio-variable-stars-fall-25/}{https://www.kaggle.com/competitions/cu-science-studio-variable-stars-fall-25/}
    \end{itemize}
  }

\newpage

\section{Достижения в школьные годы (2022--2024)}

\achievement
  {Призёр регионального этапа чемпионата «Абилимпикс»}
  {2023}
  {Компетенция: Веб-Разработка}
  {
    \begin{itemize}
      \item Третье место на региональном уровне престижного чемпионата
      \item Демонстрация навыков в современных веб-технологиях
    \end{itemize}
  }

\achievement
  {Участник всероссийской программы «СИРИУС.ИИ»}
  {2023, 2024}
  {Дважды принял участие}
  {
    \begin{itemize}
      \item Двукратное участие в престижной образовательной программе
      \item Углублённое изучение искусственного интеллекта и машинного обучения
      \item Работа над практическими проектами в области ML и AI
    \end{itemize}
  }

\achievement
  {Финалист олимпиады РТУ МИРЭА}
  {2024}
  {Математика и информатика}
  {
    \begin{itemize}
      \item Прохождение в заключительный этап олимпиады
    \end{itemize}
  }

\achievement
  {Участник зимней школы ВШЭ Санкт-Петербург}
  {2024}
  {Практическое программирование и анализ данных}
  {
    \begin{itemize}
      \item Интенсивное обучение современным методам Data Science
      \item Практика в области анализа данных и алгоритмического программирования
    \end{itemize}
  }

\achievement
  {Выпускник Московской Школы Программистов}
  {2022 -- 2023}
  {Годичный курс по Python}
  {
    \begin{itemize}
      \item Успешное прохождение профессионального курса
      \item Освоение продвинутых техник программирования и алгоритмов
    \end{itemize}
  }

\achievement
  {Участник олимпиады SMARTCUP}
  {2023}
  {МПГУ и Skysmart}
  {
    \begin{itemize}
      \item Разработка \hlteal{чат-бота для ВКонтакте} с интеграцией ИИ
      \item Работа с API социальных сетей и современными технологиями
      \item Применение искусственного интеллекта в real-time приложениях
    \end{itemize}
  }

\newpage

\section{Технические навыки}

\begin{multicols}{2}
\subsection*{Программирование}
\begin{itemize}
  \item Python (продвинутый уровень)
  \item ООП и функциональное программирование
  \item Алгоритмы и структуры данных
  \item Git/GitHub
  \item C++/Qt (нативные приложения)
\end{itemize}

\subsection*{Machine Learning}
\begin{itemize}
  \item Scikit-learn
  \item Feature Engineering
  \item Ensemble Methods
  \item Model Training \& Optimization
  \item CatBoost, XGBoost
\end{itemize}

\columnbreak

\subsection*{Data Analysis}
\begin{itemize}
  \item NumPy, Pandas
  \item Статистический анализ
  \item Временные ряды
  \item Matplotlib (визуализация)
\end{itemize}

\subsection*{NLP \& AI}
\begin{itemize}
  \item NER (CRF, SpaCy, BIO-разметка)
  \item Entity extraction / Query Understanding
  \item SpellFix, словари, regex-каскады
  \item GLiNER, sentence-transformers
  \item Prompt Engineering, работа с LLM
  \item Текстовый анализ, MLFE методология
\end{itemize}
\end{multicols}

\subsection*{Backend \& сервисы}
\begin{itemize}
  \item FastAPI, CLI, asyncio
  \item REST API, pytest
\end{itemize}

\subsection*{Математика}
\begin{itemize}
  \item Линейная алгебра и математический анализ (продвинутый уровень)
  \item Статистика и теория вероятностей (продвинутый уровень)
\end{itemize}

\subsection*{Инструменты}
\begin{itemize}
  \item Jupyter Notebook, VS Code, LaTeX, Overleaf
  \item Основы веб-разработки, работа с API, разработка чат-ботов
\end{itemize}

\newpage

\section{Профессиональные качества}

\begin{itemize}
  \item \hlblue{Усидчивость и концентрация:} способен полностью погружаться в интересующие задачи, посвящая им до 16 часов в сутки для достижения глубокого понимания
  \item \hlgreen{Самостоятельное обучение:} быстро осваиваю новые технологии и методы, активно изучаю актуальные исследования в области ML/AI
  \item \hlpurple{Командная работа:} опыт успешного взаимодействия в команде на хакатонах и олимпиадах
  \item \hlteal{Аналитическое мышление:} способность решать сложные задачи, разбивая их на подзадачи и применяя системный подход
  \item \hlorange{Стремление к развитию:} постоянно работаю над совершенствованием навыков и преодолением слабых сторон
\end{itemize}

\section{Карьерные цели}

Стремлюсь развиваться в области \hlblue{искусственного интеллекта} и \hlgreen{машинного обучения}. Заинтересован в работе в ведущих технологических компаниях России (\hlpurple{Яндекс}, \hlteal{Т-Банк}, \hlorange{Сбер}), где смогу:

\begin{itemize}
  \item Применять свои знания в ML/AI для решения реальных бизнес-задач
  \item Работать с высококвалифицированными специалистами и обмениваться опытом
  \item Участвовать в создании инновационных продуктов и исследованиях
  \item Развивать свои навыки в области Data Science и AI-систем
  \item Вносить значимый вклад в развитие российских технологий
\end{itemize}

Особый интерес представляют позиции, связанные с \hlrose{исследованиями в области ML/AI}, разработкой интеллектуальных систем, анализом данных и работой с большими языковыми моделями.

\section{Языки}

\begin{itemize}
  \item \hlblue{Русский:} родной язык
  \item \hlgreen{Английский:} технический английский (уверенное чтение документации, научных статей, участие в международных проектах)
\end{itemize}

\end{document}
